\documentclass[aspectratio=169, t, 8pt]{beamer}

%------configurando acentos-----------
\usepackage[utf8]{inputenc}
\usepackage[T1]{fontenc}
\usepackage[brazil]{babel}
\usepackage{montserrat}
\usefonttheme[onlymath]{serif}
\usepackage{booktabs} % Tabelas elegantes sem linhas verticais
\usepackage{multirow}
\usepackage{array}
\usepackage[most]{tcolorbox}
\usepackage{graphicx}
% \usepackage[inline]{enumitem}

\DeclareMathOperator{\sen}{sen}
\DeclareMathOperator{\coss}{cos}
\DeclareMathOperator{\tg}{tg}
%-------Estilos Beamer---------------
\usetheme{metropolis}
%\usecolortheme{crane}
%\usefonttheme{serif}
\linespread{1.2}

% Personalização de Cores (Exemplo: Minimal Teal)
% Cores Base
\definecolor{BgCanvas}{HTML}{F9F9FB}      % Off-white muito suave
\definecolor{MainText}{HTML}{2B2D42}      % Grafite escuro para alta legibilidade
\definecolor{PrimaryAccent}{HTML}{4A6FA5} % Azul poeira (Dusty Blue)

% Tons Pastéis para Destaques e Cards
\definecolor{PastelGreenBg}{HTML}{E8F0EC}  % Sálvia Claro (Sucesso / Destaque Positivo)
\definecolor{PastelGreenFg}{HTML}{2D5A44}

\definecolor{PastelCoralBg}{HTML}{FDEAE6}  % Coral Soft (Atenção / Problema / Alerta)
\definecolor{PastelCoralFg}{HTML}{7A3025}

\definecolor{PastelPurpleBg}{HTML}{F0EBF4} % Lavanda (Conceitos / Citações)
\definecolor{PastelPurpleFg}{HTML}{4A3E56}

\definecolor{PastelYellowBg}{HTML}{FFF4E6} % Areia / Manteiga (Notas / Dados Secundários)
\definecolor{PastelYellowFg}{HTML}{7A522A}

% --- CONFIGURAÇÕES DE TEMA BEAMER ---
%\usetheme{metropolis}

\setbeamertemplate{itemize item}{\textcolor{MainText}{$\bullet$}}
\setbeamertemplate{itemize subitem}{\textcolor{MainText}{\scriptsize$\circ$}}
\setbeamertemplate{itemize 3rditem}{\textcolor{MainText}{\tiny$\blacksquare$}}

\setbeamercolor{background canvas}{bg=BgCanvas}
\setbeamercolor{normal text}{fg=MainText}
\setbeamercolor{frametitle}{bg=PastelGreenFg, fg=BgCanvas}
\setbeamercolor{title}{fg=MainText}
\setbeamercolor{subtitle}{fg=PrimaryAccent}
\setbeamercolor{structure}{fg=PrimaryAccent}

% Remove a linha amarela padrão do Metropolis sob o título
\setbeamertemplate{title separator}{}
\setbeamertemplate{progress bar in section page}{}

% ==============================================================================
% PALETA DE CORES "GEOMETRIA AQUECIDA" (Inspirada na imagem)
% ==============================================================================
\definecolor{BgCream}{HTML}{FFFDFB}       % Fundo creme suave
\definecolor{TextDark}{HTML}{2C1A14}      % Texto em castanho escuro elegante
\definecolor{PrimaryOrange}{HTML}{E0531C}  % Laranja vibrante (Títulos e destaques)
\definecolor{Terracotta}{HTML}{B8380D}     % Terracota queimado (Alertas)
\definecolor{SoftOrange}{HTML}{FFE6D5}     % Laranja bem suave (Fundo de blocos)

% =================================================
% CONFIGURAÇÃO DO TEMA METROPOLIS
% =================================================
% Cores do texto e fundo geral
\setbeamercolor{normal text}{fg=TextDark, bg=BgCream}
\setbeamercolor{alerted text}{fg=Terracotta}
\setbeamercolor{example text}{fg=PrimaryOrange}

% Barra de título dos slides (Frametitle)
\setbeamercolor{frametitle}{fg=BgCream, bg=PrimaryOrange}

% Título principal da apresentação (Slide de Capa)
\setbeamercolor{titlelike}{fg=PrimaryOrange}
\setbeamercolor{subtitle}{fg=Terracotta}

% Barra de progresso (Linha do Metropolis)
\setbeamercolor{progress bar}{fg=PrimaryOrange, bg=PrimaryOrange!20}

% Customização de Blocos (block, alertblock, exampleblock)
\setbeamercolor{block title}{fg=BgCream, bg=PrimaryOrange}
\setbeamercolor{block body}{fg=TextDark, bg=SoftOrange}

\setbeamercolor{block title alerted}{fg=BgCream, bg=Terracotta}
\setbeamercolor{block body alerted}{fg=TextDark, bg=Terracotta!15}


% --- ESTILOS DE CARDS PASTEL COM TCOLORBOX ---
\newtcolorbox{cardPrincipal}[1][]{
  colback=SoftOrange,
  colframe=Terracotta,
  coltext=TextDark,
  arc=6pt,
  boxrule=0pt,
  left=12pt, right=12pt, top=10pt, bottom=10pt,
  fonttitle=\bfseries\small,
  #1
}

\newtcolorbox{cardcoral}[1][]{
  colback=PastelCoralBg,
  colframe=PastelCoralFg,
  coltext=PastelCoralFg,
  arc=6pt,
  boxrule=0pt,
  left=12pt, right=12pt, top=10pt, bottom=10pt,
  fonttitle=\bfseries\small,
  #1
}

\newtcolorbox{cardpurple}[1][]{
  colback=PastelPurpleBg,
  colframe=PastelPurpleFg,
  coltext=PastelPurpleFg,
  arc=6pt,
  boxrule=0pt,
  left=12pt, right=12pt, top=10pt, bottom=10pt,
  fonttitle=\bfseries\small,
  #1
}

\newtcolorbox{cardyellow}[1][]{
  colback=PastelYellowBg,
  colframe=PastelYellowFg,
  coltext=PastelYellowFg,
  arc=6pt,
  boxrule=0pt,
  left=12pt, right=12pt, top=10pt, bottom=10pt,
  fonttitle=\bfseries\small,
  #1
}

\graphicspath{{../../src/img/semana07/}}

%--------Dados do Documento
\title{Geometria ESpacial: Prismas e Pirâmides}
\subtitle{Capacidade, superfícies e a lógica do ENEM}
\author{prof. Eduardo}
\institute{Pré-Universitário Oficina do Saber / UFF}

\date{\today}


%-------- Pacotes Matemática -------------
\usepackage{amsmath}
\usepackage{amsthm}
\usepackage{amsfonts}
\usepackage{amssymb}
\usepackage{xfrac}
\usepackage{thmtools}

\begin{document}

    \begin{frame} %capa
        \vspace*{\fill}
            
            \begin{columns}[c]
                \begin{column}{0.5\textwidth}
                    {\usebeamerfont{title}\usebeamercolor[fg]{title}\inserttitle}\vspace{0.2cm}
                    
                    {\usebeamerfont{subtitle}\usebeamercolor[fg]{subtitle}\insertsubtitle}\vspace{0.5cm}

                {\usebeamerfont{author}\insertauthor}\vspace{0.3cm}

                {\usebeamerfont{institute}\insertinstitute}\vspace{0.3cm}

                {\usebeamerfont{date}\insertdate}
            \end{column}
            
            
            \begin{column}{0.4\textwidth}
                \centering
                
\includegraphics[width=0.4\textwidth]{../logo_ofSaber.png}

                \vspace{.4cm}
                \vspace*{\fill}
        \begin{columns}[c]
            \begin{column}{.5\textwidth}
                \centering
\includegraphics[width=.5\textwidth]{../Proex02c.png}
            \end{column}
            \begin{column}{.5\textwidth}
                \centering
\includegraphics[width=0.4\textwidth]{../Logo_UFF.png}
            \end{column}
        \end{columns}
        \vspace*{\fill}
            \end{column}
        \end{columns}
    \vspace*{\fill}
    \end{frame}
    
        % {
        %     \usebackgroundtemplate{
        %     \includegraphics[width=\paperwidth]{segunda_capa_pitagoras.png}}
        %     \begin{frame}[plain]                
        %     \end{frame}
        % }

    % \begin{frame}{Triângulo Retângulo: Quem é quem?}
    %     \vspace{1.3em}
    %     \begin{columns}[T, totalwidth=.9\textwidth]
    %         \begin{column}{.35\textwidth}
    %             \begin{figure}
    %                 \centering
    %                 \includegraphics[width=.92\textwidth]{slide_catetos_hipotenusa.png}
    %             \end{figure}
    %         \end{column}
    %         \begin{column}{.5\textwidth}
    %             Na aula anterior, vimos que as relações de semelhança entre triângulos permitem definir razões entre os lados de um triângulo retângulo, que só dependem dos seus ângulos agudos.

    %             \vspace{0.7em}

    %             Essas relações fundamentam a Trigonometria: uma ferramente muito útil para determinamos distâncias, inclinações, alturas, etc. Além disso, ela ainda nos permitirá estudar alguns fenômenos da natureza.

    %             \vspace{0.7em}

    %             O primeiro passo para compreendermos a trigonometria como uma importante ferramente de \textit{medida} é identificarmos com segurança cada elemento do triângulo retângulo.

    %             \vspace{0.7em}

    %             Em qualquer triângulo retângulo chamamos de \textbf{hipotenusa} ao lado oposto ao ângulo reto e aos outros dois lados, chamamos \textbf{catetos}.
    %         \end{column}
    %     \end{columns}    
    % \end{frame}

    % \begin{frame} {Posição dos Catetos}
    %     \vspace{1em}
    %     \begin{columns}
    %         \begin{column}{.35\textwidth}
    %             \begin{figure}
    %                 \centering
    %                 \includegraphics[width=\textwidth]{slide_catetos_opostosxadjacentes.png}
    %             \end{figure}
    %         \end{column}
    %         \begin{column}{.5\textwidth}
    %             Outra questão importante é a posição dos catetos, em relação ao ângulo que tomarmos como referência.

    %             \vspace{1em}

    %             Sabemos que o lado oposto ao um ângulo é o lado que o "encara". Os lados que formam o ângulo são chamados \textbf{adjacentes} a ele.

    %             \vspace{1em}
        
    %             A figura ao lado representa o cateto adjacente a cada ângulo agudo do triângulo retângulo, com uma cor semelhante. As relações de razão entre esses lados definirão as relações de seno, cosseno, tangente e suas inversas.

    %             \vspace{1em}

    %             Obviamente cada cateto possui um cateto oposto a si. Esse \textit{cateto oposto} será representado por uma cor diferente da cor do ângulo de referência.
    %         \end{column}
    %     \end{columns}
    % \end{frame}

    % \begin{frame} {Razões Trigonométricas Fundamentais}
    %     \vspace{1em}
    %     \begin{columns}[totalwidth=.6\textwidth]
    %         \begin{tabular}{m{.3\textwidth} | m{.3\textwidth} | m{.3\textwidth}}
                
    %             \textbf{Razão Trigonométrica} & \textbf{Fórmula} & \textbf{\textit{Quando usar?}}\\[1.2em]
    %             \hline
    %             \textbf{Seno} & $\sen{\left( \theta \right)} = \frac{CO}{H}$ & Relaciona o cateto adjacente de um ângulo com a hipotenusa\\[1.4em]
    %             \hline
    %             \textbf{Cosseno} & $cos{\left( \theta \right)} = \frac{CA}{H}$ & Relaciona o cateto adjacente de um ângulo com a hipotenusa\\[1.4em]
    %             \hline
    %             \textbf{Tangente} & $ \tg{\left( \theta \right)} = \frac{O}{A} $ & Relaciona o cateto oposto ao cateto adjacente de um ângulo\\[1.4em]
    %             \hline
    %         \end{tabular}
    %     \end{columns}

    %     \begin{columns}[c]
    %         \begin{column}{.6\textwidth}
    %             \begin{cardPrincipal}[title=Dica para Memorizar]
    %             \centering
    %             \textbf{\large{SOH-CAH-TOA}}

    %             \vspace{1em}
                
    %             \large{\textit{S} = $\displaystyle \frac {\text{O}} {\text{H}}$, \textit{C} = $\displaystyle \frac {\text{A}} {\text{H}}$, \textit{T} = $\displaystyle \frac {\text{O}} {\text{A}}$}
    %         \end{cardPrincipal}
    %         \end{column}
    %     \end{columns}
    % \end{frame}

    % \begin{frame} {Ângulos Notáveis}
    %     \vspace{1em}
    %     \begin{columns}[totalwidth=.9\textwidth]
    %         \begin{column}{\textwidth}
    %             \begin{tabular}{m{.25\textwidth} | m{.2\textwidth} | m{.2\textwidth} | m{.2\textwidth}}

    %                 \hline
    %                 \begin{center}
    %                     \textbf{Razão Trigonométrica}
    %                 \end{center}    
                    
    %                  & 
    %                  \begin{center}
    %                     \textbf{$30^\circ$}
    %                  \end{center}
                     
    %                  &
    %                  \begin{center}
    %                     \textbf{$45^\circ$}
    %                  \end{center}
                     
    %                  &
                     
    %                  {
    %                     \begin{center}
    %                         \textbf{$60^\circ$}
    %                     \end{center}
    %                  }
                     
    %                  \\
                    
    %                 \hline

    %                 \centering\textbf{Seno} &
    %                 \begin{center}
    %                     $\displaystyle\frac{1}{2}$    
    %                 \end{center}
    %                 &
    %                 \begin{center}
                        
    %                     $\displaystyle\frac{\sqrt{2}}{2}$
    %                 \end{center}
                    
    %                 & {
    %                     \begin{center}
    %                         $\displaystyle\frac{\sqrt{3}}{2}$
    %                     \end{center}
                        
    %                     }\\[1.8em]
                    
    %                 \hline

    %                 \begin{center}
    %                     \textbf{Cosseno}    
    %                 \end{center}
    %                 &
    %                 \begin{center}
    %                     $\displaystyle \frac{\sqrt{3}}{2}$     
    %                 \end{center}
    %                 &
    %                 \begin{center}
    %                     $\displaystyle\frac{\sqrt{2}}{2}$
    %                 \end{center} 
    %                 & {
    %                     \begin{center}
    %                         $\displaystyle \frac{1}{2}$
    %                     \end{center}
    %                 }\\[1.8em]

    %                 \hline

    %                 \begin{center}
    %                     \textbf{Tangente}
    %                 \end{center}
    %                 &
    %                 \begin{center}
    %                     $\displaystyle \frac{\sqrt{3}}{3}$
    %                 \end{center}
    %                 &
    %                 \begin{center}
    %                     $1$
    %                 \end{center}
    %                 &
    %                 {
    %                     \begin{center}
    %                         $\displaystyle {\sqrt{3}}$
    %                     \end{center}
    %                 }\\[1.8em]

    %                 \hline
    %             \end{tabular}
    %         \end{column}
    %     \end{columns}
    %     \begin{columns}[totalwidth=.9\textwidth]
    %         \begin{column}{\textwidth}
    %             % \vspace{1em}
    %             \begin{cardPrincipal}
    %                 [title=Observação]
    %                 As bancas de concurso consideram que os valores dos ângulos notáveis devem ser conhecidos pelos candidatos. Para outros ângulos, elas oferecerão valores aproximados.
    %             \end{cardPrincipal}
    %         \end{column}
    %     \end{columns}
    % \end{frame}

    % \begin{frame} {Calculando Ângulos Notáveis}
    %     \begin{columns}[T,totalwidth=.9\textwidth]
    %         \begin{column}{.45\textwidth}
    %             \centering
    %             \includegraphics[width=.9\textwidth]{slide_triangulo_equilatero.png}
    %         \end{column}
            
    %         \begin{column}{.45\textwidth}
    %             \begin{figure}
    %                 \centering
    %                 \includegraphics[width=.8\textwidth]{slide_triangulo_isosceles.png}
    %             \end{figure}
    %         \end{column}
    %     \end{columns}
    % \end{frame}

    % \begin{frame} {A Identidade Trigonométrica Fundamental}
    %     \vspace{1.2em}

    %     A identidade trigonométrica fundamental se relaciona com a base da trigonometria, que é o Teorema de Pitágoras.

    %     \begin{columns}[totalwidth=.9\textwidth]
    %         \begin{column}{\textwidth}
    %             {\Huge
    %                 \[
    %                     \left( \sen{\theta} \right)^2 + \left(\cos{\theta} \right)^2 = 1
    %                 \]
    %             }
    %         \end{column}
    %     \end{columns}
    %     \vspace{2em}
    %             Esta identidade surge naturalmente a partir da famosa relação:
    %     \begin{columns}[totalwidth=.9\textwidth]
    %         \begin{column}{\textwidth}
    %             {\Huge
    %                 \[
    %                     \left( \text{CO} \right)^2 + \left( \text{CA} \right)^2 = \text{H}^2
    %                 \]
    %             }
    %         \end{column}
    %     \end{columns}
    % \end{frame}

    % \begin{frame} {Problemas Clássicos}
    %     \begin{columns}[totalwidth=.9\textwidth]
    %         \begin{column}{\textwidth}
    %             O clichê em problemas envolvendo razões trigonométricas é a análise de situações-problemas que envolvam a medida de distâncias ou inclinações.
    %             \vspace{1em}
    %             A idéia é semapre a mesma: reduzir a situação a um trigulo retângulo, identificar os ângulos, os catetos e a hipotenusa deste triângulo e aplicar a razão trigonomêtrica mais conveniente. Vejamos um exemplo:
    %         \end{column}
    %     \end{columns}
    %     \vspace{1.2em}
    %     \begin{columns}[T, totalwidth=.9\textwidth]
    %             \begin{column}{.5\textwidth}
    %                 % \raggedleft
    %                 \includegraphics[width=.9 \textwidth]{slide_exemplo_calculoDistancia.jpeg}
    %             \end{column}
    %             \begin{column}{.45\textwidth}
    %                 Imagine que na situação representada ao lado, tenhamos que calcular a medida de um cabo que se estende, perfeitamente tensionado, de um prédio a outro.
    %                 \vspace{1em}
    %                 Que estratégia você poderia utilizar para calcular a medida deste cabo?
    %             \end{column}
    %         \end{columns}
    % \end{frame}

    % \begin{frame}{Razões Trigonométricas Recíprocas}
    %     \begin{columns}[totalwidth=.9\textwidth]   
    %         \begin{column}{\textwidth}
    %             Para finalizarmos, considero importante mencionar a existência de outras razões trigonométricas, as chamadas \textbf{razões recíprocas}. Essas são as razões que invertem a ordem fracionária: isto é, o denominador vai ao numerador e vice versa.
    %         \end{column}
    %     \end{columns}

    %     \vspace{1.4em}

    %     \centering
    %     \begin{columns}[c, totalwidth=.8\textwidth]
    %         \begin{column}{\textwidth}
    %             \begin{tabular}{m{.2\textwidth} | m{.2\textwidth} | m{.2\textwidth} | m{.2\textwidth}}
    %                 \hline
    %                 {
    %                     \begin{center}
    %                         \textbf{Razão Trigonométrica}
    %                     \end{center}
    %                 } & {
    %                     \begin{center}
    %                         \textbf{Definição}
    %                     \end{center}
    %                 } & {
    %                     \begin{center}
    %                         \textbf{Recíproca}
    %                     \end{center}
    %                 } & {
    %                     \begin{center}
    %                         \textbf{Definição}
    %                     \end{center}
    %                 }\\

    %                 \hline

    %                 {
    %                     \begin{center}
    %                         \textbf{Seno}
    %                     \end{center}
    %                 } &
    %                 {
    %                     \begin{center}
    %                         $$
    %                         \frac{
    %                             \text{Cateto Oposto}
    %                         }{
    %                             \text{Hipotenusa}
    %                         }
    %                         $$
    %                     \end{center}
    %                 } &

    %                 {
    %                     \begin{center}
    %                         \textbf{Cossecante}
    %                     \end{center}
    %                 } &
    %                 {
    %                     \begin{center}
    %                         $$
    %                         \frac{
    %                             \text{Hipotenusa}
    %                         }{
    %                             \text{Cateto Oposto}
    %                         }
    %                         $$
    %                     \end{center}
    %                 }\\

    %                 \hline

    %                 {
    %                     \begin{center}
    %                         \textbf{Cosseno}
    %                     \end{center}
    %                 } &
    %                 {
    %                     \begin{center}
    %                         $$
    %                         \frac{
    %                             \text{Cateto Adjacente}
    %                         }{
    %                             \text{Hipotenusa}
    %                         }
    %                         $$
    %                     \end{center}
    %                 } &

    %                 {
    %                     \begin{center}
    %                         \textbf{Secante}
    %                     \end{center}
    %                 } &
    %                 {
    %                     \begin{center}
    %                         $$
    %                         \frac{
    %                             \text{Hipotenusa}
    %                         }{
    %                             \text{Cateto Adjacente}
    %                         }
    %                         $$
    %                     \end{center}
    %                 }\\

    %                 \hline

    %                 {
    %                     \begin{center}
    %                         \textbf{Tangente}
    %                     \end{center}
    %                 } &
    %                 {
    %                     \begin{center}
    %                         $$
    %                         \frac{
    %                             \text{Cateto Oposto}
    %                         }{
    %                             \text{Cateto Adjacente}
    %                         }
    %                         $$
    %                     \end{center}
    %                 } &

    %                 {
    %                     \begin{center}
    %                         \textbf{Cotangente}
    %                     \end{center}
    %                 } &
    %                 {
    %                     \begin{center}
    %                         $$
    %                         \frac{
    %                             \text{Cateto Adjacente}
    %                         }{
    %                             \text{Cateto Oposto}
    %                         }
    %                         $$
    %                     \end{center}
    %                 }\\
    %                 \hline
    %                 \end{tabular}    
    %         \end{column}
    %     \end{columns}
    % \end{frame}

%----------------------------------------------
%              EXERCICIOS
%----------------------------------------------

    \begin{frame} {Exercícios} %Gabarito 1-B
        \begin{columns}
            \centering
            \begin{column}{.9\textwidth} %letra a
                1.~\textbf{(ENEM-2016): } Um petroleiro possui reservatório em formato de um paralelepípedo retangular com as dimensões dadas por 60 m x 10 m de base e 10 m de altura. Com o objetivo de minimizar o impacto ambiental de um eventual vazamento, esse reservatório é subdividido em três compartimentos, A, B e C, de mesmo volume, por duas placas de aço retangulares com dimensões de 7 m de altura e 10 m de base, de modo que os compartimentos são interligados, conforme a figura. Assim, caso haja rompimento no casco do reservatório, apenas uma parte de sua carga vazará.
            
                Suponha que ocorra um desastre quando o petroleiro se encontra com sua carga máxima: ele sofre um acidente que ocasiona um furo no fundo do compartimento C. Para fins de cálculo, considere desprezíveis as espessuras das placas divisórias.
            \vspace{1.2em}
            \begin{columns}
                \centering
                \begin{column}{.55\textwidth}
                    \begin{figure}
                        \includegraphics[width=\textwidth]{questao01.png}
                    \end{figure}
                \end{column}
                \begin{column}{.45\textwidth}
                    Após o fim do vazamento, o volume de petróleo derramado terá sido de 
                    \begin{enumerate}[a.]
                        \item $1,4 \cdot 10^{3} m^{3}$
                        \item $1,8 \cdot 10^{3} m^{3}$
                        \item $2,0 \cdot 10^{3} m^{3}$
                        \item $3,2 \cdot 10^{3} m^{3}$
                        \item $6,0 \cdot 10^{3} m^{3}$
                    \end{enumerate}
                \end{column}
            \end{columns}
            \end{column}    
        \end{columns}
    \end{frame}

    \begin{frame} {Exercícios} %Gabarito 2-A
        \begin{columns}[totalwidth=.9\textwidth]
            \begin{column}{\textwidth}
                2.\textbf{(UFPR-2018) }A figura abaixo apresenta um molde para construção de uma pirâmide hexagonal regular. Para montar essa pirâmide, basta recortar o molde seguindo as linhas contínuas, dobrar corretamente nas linhas tracejadas e montar a pirâmide usando as abas trapezoidais para fixar sua estrutura com um pouco de cola. Sabendo que cada um dos triângulos tracejados nesse molde é isósceles, com lados medindo 5 cm e 13 cm, qual das alternativas abaixo mais se aproxima do volume dessa pirâmide? 
            \end{column}
        \end{columns}
        \vspace{1.6em}
        \begin{columns}[T]
            \begin{column}{.65\textwidth}
                \begin{figure}
                    \centering
                    \includegraphics[width=.5\textwidth]{questao02.png}                    
                \end{figure}
            \end{column}

            \begin{column}{.28\textwidth}
                \begin{enumerate}[a. ]
                    \item $260 \text{ cm}^3$
                    \item $276 \text{ cm}^3$
                    \item $281 \text{ cm}^3$
                    \item $390 \text{ cm}^3$
                    \item $780 \text{ cm}^3$
                \end{enumerate}
            \end{column}
        \end{columns}
    \end{frame}

    \begin{frame} {Exercícios} %Gabarito 3-C
        \begin{columns}[totalwidth=.9\textwidth]
            \begin{column}{\textwidth}
                3 \textbf{(ENEM-2012) }Um reservatório de uma cidade estava com 30 m³ de água no momento em que iniciou um vazamento estimado em 30 litros por minuto. Depois de 20 minutos, a partir do início do vazamento, uma equipe técnica chegou ao local e gastou exatamente 2 horas para consertar o sistema e parar o vazamento. O reservatório não foi reabastecido durante todo o período em que esteve com o vazamento.

                \vspace{1em}
                
                Qual foi o volume de água que sobrou no reservatório, em m³, no momento em que parou o vazamento?
            \end{column}   
        \end{columns}
            \vspace{1.8em}

        \begin{center}
            \begin{columns}
                \begin{column}{.15\textwidth}
                    \centering
                    \textbf{a. } $3,6.$
                \end{column}
                \begin{column}{.15\textwidth}
                    \centering
                    \textbf{b. } $4,2.$
                \end{column}
                \begin{column}{.15\textwidth}
                    \centering
                    \textbf{c. } $25,8.$
                \end{column}
                \begin{column}{.15\textwidth}
                    \centering
                    \textbf{d. } $26,4.$
                \end{column}
                \begin{column}{.15\textwidth}
                    \centering
                    \textbf{e. } $27,6.$
                \end{column}
            \end{columns}
        \end{center}
    \end{frame}

    \begin{frame} {Exercícios} %Gabarito 4-A
        \begin{columns}[totalwidth=.9\textwidth]
            \begin{column}{\textwidth}
                4 \textbf{(ENEM-2014) }Uma pessoa comprou um aquário em forma de um paralelepípedo retângulo reto, com $40$ cm de comprimento, $15$ cm de largura e $20$ cm de altura. Chegando em casa, colocou no aquário uma quantidade de água igual à metade de sua altura. A seguir, para enfeitá-lo, irá colocar pedrinhas coloridas, de volume igual à $50 \text{ cm}^3$ cada, que ficaração totalmente submersas no aquário.
                
                \vspace{1em}

                Após a colocação das pedrinhas, o nível da água deverá ficar a $6$ cm do topo do aquário.

                \vspace{1em}

                O número de pedrinhas a serem colocadas deve ser igual a
            \end{column}   
        \end{columns}
            \vspace{1.8em}

        \begin{center}
            \begin{columns}
                \begin{column}{.15\textwidth}
                    \centering
                    \textbf{a. } $48.$
                \end{column}
                \begin{column}{.15\textwidth}
                    \centering
                    \textbf{b. } $72.$
                \end{column}
                \begin{column}{.15\textwidth}
                    \centering
                    \textbf{c. } $84.$
                \end{column}
                \begin{column}{.15\textwidth}
                    \centering
                    \textbf{d. } $120.$
                \end{column}
                \begin{column}{.15\textwidth}
                    \centering
                    \textbf{e. } $168.$
                \end{column}
            \end{columns}
        \end{center}
    \end{frame}

    \begin{frame} {Exercícios} %Gabarito 5-B
        \begin{columns}[totalwidth=.9\textwidth]
            \begin{column}{\textwidth}
                5 \textbf{(ENEM-2010) } A siderúrgica ''Metal Nobre'' produz diversos objetos maciços utilizando o ferro. Um tipo especial de peça feita nessa companhia tem o formato de um paralelepípedo retangular, de acordo com as dimensões indicadas na figura que segue.

                O produto das três dimensões indicadas na peça resultaria na medida da grandeza
            \end{column}
        \end{columns}

        \vspace{1em}
        
        \begin{columns}[totalwidth=.9\textwidth]
            \begin{column}{.5\textwidth}
                \begin{figure}
                    \centering
                    \includegraphics[width=\textwidth]{questao05.png}
                \end{figure}
            \end{column}
            \begin{column}{.3\textwidth}
                \begin{enumerate}[a. ]
                    \item massa.
                    \item volume.
                    \item superfície.
                    \item capacidade.
                    \item comprimento.
                \end{enumerate}
            \end{column}
        \end{columns}
    \end{frame}

    \begin{frame} {Exercícios} 
        6 \textbf{(ENEM-2010) }Um porta-lápis de madeira foi construído no formato cúbico, seguindo o modelo ilustrado a seguir. O cubo de dentro é vazio. A aresta do cubo maior mede 12 cm e a do cubo menor, que é interno, mede 8 cm.
        \begin{columns}[totalwidth=\textwidth]
            \begin{column}{.35\textwidth}
                \begin{figure}
                    \raggedleft
                    \includegraphics[width=\textwidth]{questao06.png}
                \end{figure}
            \end{column}
            \begin{column}{.6\textwidth}
                O volume de madeira utilizado na confecção desse objeto foi de:
                \begin{enumerate}[a. ]
                    \item $12\text{ cm}^3$
                    \item $64\text{ cm}^3$
                    \item $96\text{ cm}^3$
                    \item $1216\text{ cm}^3$
                    \item $1728\text{ cm}^3$
                \end{enumerate}
            \end{column}
        \end{columns}
    \end{frame}
\end{document}